\section{Actual-source arithmetic and Boolean compilation}\label{sec:source-arithmetic}

\begin{definition}[Actual-source arithmetic]\label{def:source-language}
Let $\Esrc$ be existential arithmetic over $\N$ with constants, addition,
equality, order and finite Boolean combinations.  An extra function may be
given by a supplied closed original unary source $H:o\to o$:
\[
 F_H(n)=\can{HX}_0-1,\qquad x_0=n+1.
\]
We also allow an observation leaf $\operatorname{Source}(A)=\can A_0-1$
for any supplied ground-output original source with ground free names,
using the same natural/raw encoding for those names.
\end{definition}

Source equations in the constructions below are pointwise equations
at one joint valuation, rather than universal comparison assertions.
Their auxiliary raw coordinates are positive and represent natural
variables plus one. Translating a displayed raw linear equation into
that convention only adds fixed integer constants on its two sides.

The source syntax is part of the finite input.  One does not adjoin every
ordinal-indexed collapse merely because its notation is computable.  The functions are total. The compiler below is valid for each fixed
control satisfying Section~\ref{sec:setting}; it does not by itself
decide the resulting existential sentences.

\subsection{Literal arithmetic and collision gadgets}\label{sec:gadgets}

Write $g(n)=\psi_f(\omega+n)=f(n+2)+1$.  From the canonical clauses,
for a fresh unused ground binder define
\[
 e(A,B)=(\lambda u:o.A)B,\qquad
 \can{e(A,B)}_0=g(3\can A_0+\can B_0+2).
\]
This uses the full canonical assignment rather than evaluating the beta
reduct.  The exact coefficient-one formula makes $g$ injective. Our adjacent
growth assumption gives $f(t+1)>4(f(t)+1)$, and hence
$g(n+1)>4g(n)$ at every argument used below.

\begin{lemma}[Addition and small constants]\label{lem:addition}
Natural addition $n_a+n_b=n_c$, denoted $\operatorname{Plus}(a,b,c)$
on its raw representatives, is defined by the two source equations
\[
 e(b,0)=e(0,d),\qquad e(a,d)=e(c,0),
\]
with one fresh positive raw witness $d$.  Arithmetic one has a constant-size
source-equation definition.  Every binary literal $k$ has an addition DAG
of size $O(\log(k+1))$ once zero and one are available.
\end{lemma}
\begin{proof}
Injectivity gives $d=3b-2\ge1$ and $a+b=c+1$.  The latter is exactly the
natural addition equation after subtracting one from each raw coordinate.
The witness is unique, and any true addition supplies it.

Let $I=\lambda z:o.z$.  The prime and lambda clauses give
$\can S=(1,1)$ and $\can I=(3,1)$.  The equation $Sv=I0$ forces raw
$v=3$.  Adding $\operatorname{Plus}(u,u,v)$ forces raw $u=2$, representing
natural one.  Its addition witness is $d=4$.  Thus no giant closed value
has to be evaluated to define one.  Repeated doubling and add-one gates
give the usual binary construction of $k$.
\end{proof}

A supplied graph call $F_H(n_a)=n_b$ is simply the source equation $Ha=b$.
A source leaf is flattened by $A=y$.  These equations retain the actual
supplied source and every shared input occurring inside it.

\begin{lemma}[Superincreasing pair coding]\label{lem:sort}
Let $q(n)=F_S(n)=f(n+4)$.  Then
\[
 q(a)+q(b)=q(c)+q(d)
 \quad\Longleftrightarrow\quad\{a,b\}=\{c,d\}
\]
as multisets.  Consequently sorting two naturals has a primitive-positive
definition using addition and the actual $S$ graph.
\end{lemma}
\begin{proof}
The control bound gives $q(n+1)>2q(n)$.  Cancel equal indices across the
two sides.  If a largest remaining index $M\ge1$ occurs, the other side
has at most two values at most $q(M-1)$, too small to equal $q(M)$.
If only index zero remains on one side, its positive contribution also
cannot vanish.  This proves multiset equality; the converse is immediate.
Impose
\[
 q(a)+q(b)=q(l)+q(h),\qquad l+\delta=h.
\]
Multiset equality and nonnegative $\delta$ force $l=\min(a,b)$,
$h=\max(a,b)$ and $\delta=|a-b|$.  These values conversely satisfy the
equations.  Flattening the sums and graph calls uses only the previous lemma.
\end{proof}

\begin{lemma}[Six-eraser collision detector]\label{lem:collision}
For arbitrary actual ground sources $P,Q$, put
\[
 L=e(e(P,P),e(Q,Q)),\qquad R=e(e(P,Q),e(Q,P)).
\]
At every joint valuation $L_0\ge R_0$, with equality exactly when
$P_0=Q_0$.  Otherwise $L_0-R_0\ge2$.
\end{lemma}
\begin{proof}
For raw $p=P_0,q=Q_0$ set
\[
 A=g(4p+2),\ B=g(4q+2),\ C=g(3p+q+2),\ D=g(p+3q+2).
\]
Then $L_0=g(3A+B+2)$ and $R_0=g(3C+D+2)$, with all inner arguments
at least six.  If $p>q$, then $A>4C$ and $C\ge D$, giving
$3A+B>3C+D$.  If $q>p$, then $B>4D$ and $D\ge C$, giving the same
strict conclusion.  For $p=q$ every inner value agrees.  The fourfold
growth bound on the outer $g$ also gives the integer gap at least two.
Thus the unequal weights three and one do not require a symmetric source adder.
\end{proof}

\begin{lemma}[Conjunction without cancellation]\label{lem:aggregate}
A finite conjunction of actual ground-source equations can be transformed
into ground sources $L,R$ with $L_0\ge R_0$ everywhere, equality precisely
when every equation holds, and otherwise gap at least two.
\end{lemma}
\begin{proof}
Apply the collision detector to each equation.  Combine all left outputs
and all right outputs in identically shaped $e$-trees.  Strict increase
in both arguments preserves nonnegativity, and equality survives exactly
when it survives at every leaf.  One unequal leaf makes its ancestors
strict; the growth bound gives the final gap.  A balanced tree adds only
linear DAG overhead.  For an empty conjunction take $L=R=0$.
\end{proof}

\subsection{A full Boolean compiler}\label{sec:boolean}

\begin{theorem}[One unconditional comparison]\label{thm:compiler}
Every existential $\Esrc$ sentence has a computably constructed ground-output
strict comparison, and one weak comparison, true precisely when the sentence
is false.  All free variables of the resulting sources are ground.  Shared
inputs and source leaves retain their valuations.  The added source DAG size
is polynomial in the formula DAG, binary constants and supplied source DAGs.
\end{theorem}
\begin{proof}
Flatten the arithmetic terms once, sharing each original DAG node.  The
addition and supplied-source graph equations are conjoined outside the
Boolean matrix.  They define total functions, so this operation preserves
truth even under negation.

For each Boolean node construct a natural bit $b$.  To test equality $a=c$,
sort $a,c$ and obtain $\delta=|a-c|$; sort $\delta,1$ and let
$v=\min(\delta,1)$; require $b+v=1$.  Thus $b=1$ exactly when the inputs
are equal.  Weak order $a\le c$ is the equality test $\max(a,c)=c$.
Strict order is the complement of the reverse weak-order bit.  Negation
uses $b+c=1$, conjunction uses $b=\min(c,d)$, and disjunction uses
$b=\max(c,d)$.  The preceding sorting lemma defines each operation by a
constant number of equations.  Induction proves that every output is in
$\{0,1\}$ and has exactly its intended truth value.

Require the root bit to be one.  For each valuation of the original inputs,
the flattened values, sorting witnesses, differences and bits are uniquely
determined, and supply a solution exactly when the matrix is true.  Addition
witnesses are positive by Lemma~\ref{lem:addition}.  This proves both directions
without a disjunctive-normal-form expansion or an unbounded integer witness.
Repeated source names have never been split into independent copies.

Apply Lemma~\ref{lem:aggregate} to the resulting conjunction.  Since every
existential witness is now a free ground coordinate, universal $R\mslt L$
is exactly absence of a witness.  For the weak version put
$K=\lambda z:o.0$, with $\can K=(5,1)$.  Literal application gives
\[
 \can{KR}_0\le\can{IL}_0
 \quad\Longleftrightarrow\quad g(5+R_0)\le g(3+L_0)
 \quad\Longleftrightarrow\quad L_0-R_0\ge2.
\]
Equality or a gap of two are the only alternatives.  The weak comparison
therefore expresses the same unsatisfiability assertion.

Every arithmetic or Boolean node adds constant gadget overhead, binary
constants add logarithmic overhead, and the final detector/aggregation adds
linear overhead.  Existing source DAGs remain shared.  This proves the
polynomial DAG claim, not a polynomial expanded-tree or evaluation-time bound.
\end{proof}

\begin{corollary}[Ground weak/strict equivalence]\label{cor:ground-equivalence}
On all actual ground-output sources with ground free variables, universal
weak and strict comparison are polynomial-DAG many-one interreducible.
\end{corollary}
\begin{proof}
For a weak query $A\mle B$, compile the existential counterexample
$\operatorname{Source}(A)>\operatorname{Source}(B)$ into the strict version.
For a strict query compile $\operatorname{Source}(A)\ge\operatorname{Source}(B)$
into the weak version.  No recognition of a positive arithmetic fragment is
needed, because the compiler accepts arbitrary supplied ground-source leaves.
\end{proof}

